% Quelle: Gedächtnisprotokoll_klausur_oder_probeklausur_unsicher.md
% Art: Klausur/Klausurvorbereitung
% Themen: RSA, Primfaktorzerlegung, Inklusion/Exklusion, diophantische Gleichungen, ggT, lineare Codes, Kontrollmatrix, Syndrom
\section{Gedächtnisprotokoll (Klausur oder Probeklausur)}
% Gedächtnisprotokoll: Aufgaben nur sinngemäß erinnert; keine Punktzahlen und keine Lösungen angegeben.

\subsection{Aufgabe 1}
\begin{aufgabe}[Aufgabe 1]
\begin{itemize}
  \item RSA Aufgabe mit $m=2803, d=113$ und $w=715$
  \item man sollte die Formeln für die Verschlüsselung und Entschlüsselung angeben
  \begin{itemize}
    \item \textit{Notiz: einfach aus dem Skript kopieren?}
  \end{itemize}
  \item Und man sollte sagen, wie Oscar (man in the middle) $d$ bestimmen kann wenn er $m$ kennt
  \begin{itemize}
    \item man sollte zwei Lösungswege angeben
    \item \textit{Notiz: I guess die zwei cases mit $m$ ist Primzahl und $m$ ist nicht Primzahl?}
  \end{itemize}
  \item Irgendwo hier zwischen sollte man eine Primzahlzerlegung machen
\end{itemize}
% Punktzahl: nicht angegeben.
\end{aufgabe}

\subsection{Aufgabe 2}
\begin{aufgabe}[Aufgabe 2]
\begin{itemize}
  \item Inklusion / Exklusion
  \item Bestimmen Sie, welche Zahlen $\le 840$ nicht durch 2, 3 und 5 teilbar sind
  \begin{itemize}
    \item verwenden Sie die Zeichen $A_i$ und $\alpha_i$ wie in der Vorlesung gezeigt
  \end{itemize}
\end{itemize}
% Punktzahl: nicht angegeben.
\end{aufgabe}

\subsection{Aufgabe 3}
\begin{aufgabe}[Aufgabe 3]
\begin{itemize}
  \item diophantische Gleichung lösen
  \item \textit{Notiz: weiß leider nicht mehr welche das war, aber war sehr easy}
  \item war von der Form $ax-by=30$
  \item man sollte einmal in a) den $\ggT$ von $a$ und $b$ bestimmen
  \item und dann in b) eine ganzzahlige Lösung bestimmen
\end{itemize}
% UNSICHER: konkrete Werte von a und b nicht überliefert.
% Punktzahl: nicht angegeben.
\end{aufgabe}

\subsection{Aufgabe 4}
\begin{aufgabe}[Aufgabe 4]
\begin{itemize}
  \item $[10, 8, 3]_{11}$ Code und die folgenden Wörter wurden empfangen
  \begin{itemize}
    \item $r=115000711$
    \item $r=115000733$
  \end{itemize}
\end{itemize}
% UNSICHER: Empfangene Wörter haben in der Quelle nur 9 Stellen, obwohl der Code Länge 10 hat; im Text als r_1 und r_2 bezeichnet.
\begin{enumerate}
  \item[a)] Stellen Sie die Kontrollmatrix auf
  \item[b)] Berechnen Sie Syndrom von $r_1$ und $r_2$
  \item[c)] Welches der beiden Wörter ist nicht dekodierbar und wieso?
  \item[d)] Korrigieren Sie das entsprechende Wort, das dekodierbar ist (also bestimmen Sie $c$)
  \item[c)] Zeigen Sie, dass das korrigierte Wort tatsächlich ein Codewort ist
  % UNSICHER: Teilaufgabe in der Quelle erneut als "c)" bezeichnet (vermutlich e)).
\end{enumerate}
% Punktzahl: nicht angegeben.
\end{aufgabe}
